\sectionguard
\subsection{Athanasius of Alexandria}\label{sec:fb-athanasius}

After Nicaea the teaching on the angels is drawn into the defence of the
Son's Godhead. The Arians held the Word to be the first and greatest of
the creatures, the instrument through whom God made the rest. Athanasius
answers by setting the angels, with all their orders, among the things
made, and the Son on the side of the Maker. What he says of the angels is
therefore said with care: they are creatures, many and ranked, ministers
sent and worshippers of the Son; they do not create, and God needs none
of them in order to create. In the \emph{Life of Antony} the same bishop
hands on the desert's teaching on the fallen angels, their wiles, their
weakness, and the discernment of spirits. English quotations of the
Nicene Fathers are from Philip Schaff's \emph{Nicene and Post-Nicene
Fathers}, cited by the witness's own divisions;
passages given from Migne's Greek are rendered for this edition, as the
notes mark.

\subsubsection{The angels among the things made}

The Arian formula called the Son ``a creature, but not as one of the
creatures.'' Athanasius replies that the phrase distinguishes nothing,
for no creature is simply like another. The six days bring forth light
and night, sun and moon, beasts and man, none as another; and in the
unseen creation likewise, ``nor the Angels as the Thrones, nor the Thrones
as the Authorities, yet they are all creatures, but each of the things made
according to its kind exists and remains in its own essence, as it was
made'' (\emph{Orationes contra Arianos} II.19). Rank among creatures is
real, and it is not causality. ``Star, for instance, differs from star in
glory,'' yet ``it follows not for all this that some are lords, and
others servants to the superior, nor that some are efficient causes,
others by them come into being, but all have a nature which comes to be
and is created, confessing in their own selves their Framer'' (II.20).

From this Athanasius draws the teaching that the angels do not create:

\begin{quote}
For none of things which are brought to be is an efficient cause, but all
things were made through the Word: who would not have wrought all things,
were He Himself in the number of the creatures. For neither would the
Angels be able to frame, since they too are creatures, though Valentinus,
and Marcion, and Basilides think so, and you are their copyists; nor will
the sun, as being a creature, ever make what is not into what is.
(\emph{C.\ Ar.}\ II.21)
\end{quote}

\noindent The Arians answered that created nature could not bear to be
made by the unmediated hand of God, and so needed a mediator. Athanasius
turns the argument on itself. If the Word is a creature, then he too
needs a medium in his own making, and that medium another, ``and thus
tracing back and following out, we shall invent a vast crowd of
accumulating mediators; and thus it will be impossible that the creation
should subsist'' (II.26). God is neither too proud to make the least
things nor too weak to make them himself; the Lord who numbers the
sparrows and clothes the grass of the field is the one who made them
(II.25). The Arian picture would ascribe to God ``jealousy, in that He has
not taught many how to frame, so that there may be around Him, as
Archangels and Angels many, so framers many; and weakness, in that He
could not make by Himself, but needed a fellow-worker, or under-worker.''
``But God is deficient in nothing: perish the thought!'' (II.29).

The distinction that settles the matter is between making and ministry.
Moses led the people and delivered the Law, but ``to minister is of
things originate as of servants, but to frame and to create is of God
alone, and of His proper Word and His Wisdom'' (II.27). Ministers God has
in abundance: ``For there are many Archangels, many Thrones, and
Authorities, and Dominions, thousands of thousands, and myriads of
myriads, standing before Him, ministering and ready to be sent'' (II.27;
cf.\ Dan 7:10). The multitude of the angels is the sign of their
creatureliness; the Word is one, and there are ``not many Words, but one
only Word of the one Father, and one Image of the one God'' (II.27).
Aquinas determines the same point: creation is the act
proper to God alone, and no creature creates, not even as God's
minister (\ST{I q.~45 a.~5}); the Fourth Lateran Council confesses the one
God as ``creator omnium visibilium et invisibilium, spiritualium et
corporalium'' (DS~800).

\subsubsection{The Son is no angel}

The Arians pressed Heb 1:4, where the Son is ``made so much better than
the angels,'' as a comparison between beings of one kind. Athanasius
first fixes the time and person of the text: the Apostle speaks of the
Word made flesh, whose ministry surpasses the ministry of the Law ``spoken
by angels'' (Heb 2:2). Then he weighs the word itself. ``This is why
throughout he uses no comparison, such as `become greater,' or `more
honourable,' lest we should think of Him and them as one in kind, but
`better' is his word, by way of marking the difference of the Son's
nature from things originated'' (\emph{C.\ Ar.}\ I.55). Scripture says
wisdom is ``better'' than rubies and one day in God's courts ``better''
than a thousand, though stones and wisdom share no nature:

\begin{quote}
In like manner there is nought akin between the Son and the Angels; so
that the word `better' is not used to compare but to contrast, because of
the difference of His nature from them. And therefore the Apostle also
himself, when he interprets the word `better,' places its force in
nothing short of the Son's excellence over things originated, calling the
one Son, the other servants; the one, as a Son with the Father, sitting
on the right; and the others, as servants, standing before Him, and being
sent, and fulfilling offices. (I.55)
\end{quote}

\noindent Those who make the Son one of the angels in nature repeat the
Gnostics, for Valentinus ``said that the Angels were one in kind with the
Christ, and Carpocrates that Angels are framers of the world'' (I.56).
The whole contrast of Heb 1 is between the one who sits and those who
stand: ``If the Son be in the number of the Angels, then let the word
`become' apply to Him as to them \dots\ or be the Son standing with them
all as a ministering Spirit, sent forth to minister Himself as they
are'' (I.62). The angels minister ``ascending and descending,'' and when
they come they say ``I am sent unto thee''; of the Son it is written,
``And let all the Angels of God worship Him'' (I.61; Heb 1:6).

Worship is the decisive sign. Had the Son been only the highest
creature, Scripture would have called him ``greater than Archangels, and
more honourable than the Thrones''; instead it calls him Son, and ``the
Angels ministered unto Him, as being one beyond themselves; and they
worship Him, not as being greater in glory, but as being some one beyond
all the creatures, and beyond themselves, and alone the Father's proper
Son according to essence'' (II.23). Creature does not worship creature,
``but servant Lord, and creature God.'' The angel of the Apocalypse
forbids John to worship him, and the angel of Manoah refuses the
sacrifice: ``Therefore to God alone appertains worship, and this the
very Angels know, that though they excel other beings in glory, yet they
are all creatures and not to be worshipped, but worship the Lord''
(II.23; Apoc 22:9; Judg 13:16).

The same reasoning governs the angels of the patriarchs. When Jacob
blesses the sons of Joseph, ``the angel that delivereth me from all
evils, bless these boys'' (Gen 48:16), Athanasius holds that ``none of
created and natural Angels did he join to God their Creator''; ``he
shewed that it was no created Angel, but the Word of God,'' the Angel of
great Counsel, ``whom he joined to the Father in his prayer'' (\emph{C.\
Ar.}\ III.12). No one prays to receive a gift from God and an angel
together, for ``it is proper then to an Angel to minister at the command
of God, and often does he go forth to cast out the Amorite, and is sent
to guard the people in the way; but these are not his doings, but of God
who commanded and sent him'' (III.12). When the angels appear, the
seer knows that he has seen an angel and not God; Zacharias saw an angel
and Isaiah saw the Lord. At the bush ``the Angel of the Lord was seen in a
flame of fire,'' and the Lord called to Moses; ``yet was not the Angel
the God of Abraham, but in the Angel God spoke. And what was seen was an
Angel; but God spoke in him'' (III.14; Exod 3:2--6). The angels, ``as it
is written, are `ministering spirits sent forth to minister,' and are
heralds of gifts given by Him through the Word to those who receive
them'' (III.14; Heb 1:14). Aquinas too refers the apparitions to the Son:
``all the apparitions in the Old Testament were ordained to that one
whereby the Son of God appeared in the flesh'' (\ST{I q.~51 a.~2 ad~1}).

\subsubsection{The fall and the air}

The good angels keep their place by willing what God wills. Athanasius
says so in passing, against the Arian claim that the Son is one with the
Father only by agreement of will; were that the measure, the Thrones and
Dominions would be sons too, for ``what God wills, that will they \dots\
For they would not have remained in their own glory, unless, what the
Father willed, that they had willed also.'' The one who did not remain
heard the word of Isaiah against Lucifer, fallen from heaven (\emph{C.\
Ar.}\ III.10; Isa 14:12). The fallen spirits now hold the lower air.
Explaining why the Lord chose death on the Cross, Athanasius writes:

\begin{quote}
And once more, if the devil, the enemy of our race, having fallen from
heaven, wanders about our lower atmosphere, and there bearing rule over
his fellow-spirits, as his peers in disobedience, not only works
illusions by their means in them that are deceived, but tries to hinder
them that are going up \dots\ while the Lord came to cast down the
devil, and clear the air and prepare the way for us up into heaven
\dots\ by what other kind of death could this have come to pass, than by
one which took place in the air, I mean the cross? (\emph{De
incarnatione} 25.5)
\end{quote}

\noindent ``For thus being lifted up He cleared the air of the malignity
both of the devil and of demons of all kinds'' (25.6; Eph 2:2; Luke
10:18). Aquinas assigns the demons the same double abode: hell by reason
of their sin, and ``the darksome atmosphere'' for the exercise of men
until the judgment (\ST{I q.~64 a.~4}).

\subsubsection{Antony's discourse on the demons}

In the long discourse that Athanasius places in the mouth of Antony
(\emph{Vita Antonii} 16--43) the father of monks teaches his disciples
about the fallen angels. It is spoken to monks, and its aim is practical:
not to describe the demons' nature, ``for others of greater powers than we
possess,'' but ``to know their wiles against ourselves'' (21). Its
teaching may be gathered under four heads.

\paragraph{Their origin and envy.} ``Great is their number in the air
around us, and they are not far from us'' (21; Eph 6:12). Antony begins
with the faith of the Church on their origin:

\begin{quote}
First, therefore, we must know this: that the demons have not been
created like what we mean when we call them by that name; for God made
nothing evil, but even they have been made good. Having fallen, however,
from the heavenly wisdom, since then they have been grovelling on earth.
(22)
\end{quote}

\noindent Their motive is envy: ``out of envy of us Christians they move
all things in their desire to hinder us from entry into the heavens; in
order that we should not ascend up thither from whence they fell'' (22).
Their first weapon is the evil thought; when that fails, the phantasm,
``taking the forms of women, wild beasts, creeping things, gigantic
bodies, and troops of soldiers.'' ``For they are nothing and quickly
disappear, especially if a man fortify himself beforehand with faith and
the sign of the cross'' (23). They imitate psalmody and echo the reading
of Scripture, rouse the monk to prayer, urge immoderate fasting, and put
on the appearance of holy men, ``not for the sake of piety or truth, but
that they may carry off the simple to despair'' (25). Even when they
speak truly, as when they confessed Christ the Son of God, the Lord
silenced them, ``that He might accustom us never to give heed to them
even though they appear to speak what is true'' (26).

\paragraph{Their weakness.} The centre of the discourse is the weakness
of the demons since Christ. ``Since the Lord visited earth, the enemy is
fallen and his powers weakened'' (28). Antony argues from their own
behaviour. Spirits who pass through shut doors and haunt the whole air,
and who hate the servants of God, would destroy them at once if they
could; that they threaten, gather in crowds, and change their shapes
shows that they cannot:

\begin{quote}
But the demons as they have no power are like actors on the stage
changing their shape and frightening children with tumultuous apparition
and various forms: from which they ought rather to be despised as shewing
their weakness. At least the true angel of the Lord sent against the
Assyrian had no need for tumults nor displays from without, nor noises
nor rattlings, but in quiet he used his power and forthwith destroyed a
hundred and eighty-five thousand. (28; 4 Kgs 19:35)
\end{quote}

\noindent The history of Job teaches the same: ``the devil was not the
strong man, but God who delivered Job to him to be tried,'' and he had to
ask permission twice. The demons of Gerasa could not enter the swine
without leave; ``if they had power not even against swine, much less have
they any over men formed in the image of God'' (29). The Saviour drew
the devil ``like a dragon'' ``with a hook,'' and he was bound by the Lord ``as a sparrow, that
we should mock him'' (24). The fire in which the demons appear is ``no
true light,'' but ``the preludes and likenesses of the fire prepared for
the demons'' (24). What they fear is the life of the servants of Christ:
``the fasting, the sleeplessness, the prayers, the meekness, the
quietness, the contempt of money and vainglory, the humility, the love of
the poor, the alms, the freedom from anger of the ascetics, and, chief of
all, their piety towards Christ'' (30). ``For they are cowards, and
greatly fear the sign of the Lord's Cross, since of a truth in it the
Saviour stripped them, and made an example of them'' (35; Col 2:15).

Antony confirms the teaching from what he himself endured. A demon
``exceeding high'' announced, ``I am the power of God and I am
Providence,'' and at the name of Christ ``he, big as he was, together
with all his demons, disappeared'' (40). Satan himself came to his door
complaining that the monks cursed him undeservedly: ``I am not he who
troubles them, but they trouble themselves, for I am become weak \dots\
The Christians are spread everywhere, and at length even the desert is
filled with monks.'' Antony replied, ``For the coming of Christ hath made
thee weak, and He hath cast thee down and stripped thee''; and the devil,
``having heard the Saviour's name, and not being able to bear the burning
from it, vanished'' (41). The enemy adapts himself to the soul he finds:
``they approach us in a form corresponding to the state in which they
discover us,'' so that the fearful are terrified the more and the joyful
put them to flight. ``Thus the enemy, seeing Job fenced round with them,
withdrew from him; but finding Judas unguarded, him he took captive''
(42).

\paragraph{Their foreknowledge.} The demons' predictions deceive the
simple, and Antony explains them without granting the demons any power
over the future. ``For they know none of those things which are not yet in
existence; but God only is He who knoweth all things before their
birth'' (31). They see a traveller set out and outrun him; they see rain
in Ethiopia and announce the Nile's flood before it reaches Egypt (31--32).
When they guess aright they do what physicians, pilots, and farmers do:
``For experienced physicians also, since they see the same malady in
different people, often foretell what it is, making it out by their
acquaintance with it'' (33). Aquinas determines the same with the same
image: angels know the future in its causes, surely where the causes are
necessary and conjecturally where they usually hold, ``thus the doctor
knows beforehand the health of the patient''; the future in itself
belongs to God alone (\ST{I q.~57 a.~3}). A soul ``perfectly pure and in
its natural state'' sees further than the demons, ``for it has the Lord
who reveals to it,'' as Elisha saw the hosts that stood about him
(\emph{Vita Antonii} 34; 4 Kgs 6:17).

\paragraph{The discernment of visitations.} Antony's teaching on
discernment distinguishes the visit of a holy angel from the assault of a
demon. When spirits come by night and say, ``we are the angels,'' the
monk is to sign himself and pray, and they vanish (35). The two
visitations differ in their manner and in their fruit:

\begin{quote}
The vision of the holy ones is not fraught with distraction \dots\ But it
comes so quietly and gently that immediately joy, gladness and courage
arise in the soul. For the Lord who is our joy is with them, and the
power of God the Father. (35)

But the inroad and the display of the evil spirits is fraught with
confusion, with din, with sounds and cryings such as the disturbance of
boorish youths or robbers would occasion. From which arise fear in the
heart, tumult and confusion of thought, dejection, hatred towards them
who live a life of discipline, indifference, grief, remembrance of
kinsfolk and fear of death, and finally desire of evil things, disregard
of virtue and unsettled habits. (36)
\end{quote}

\noindent Fear may attend both, but not in the same way. When Gabriel
came to Zacharias and to Mary, and the angel to the women at the tomb,
their fear ``arose not from timidity, but from the recognition of the
presence of superior beings,'' and the angels at once took it away (35).
``For joy and a settled state of soul show the holiness of him who is
present'' (36). The demons never remove fear; they increase it, and at
last demand, ``fall down and worship,'' as they deceived the Greeks into
worshipping them (37). Hence the sign: ``whenever there is any
apparition, be not prostrate with fear, but whatsoever it be, first
boldly ask, Who art thou? And from whence comest thou? And if it should
be a vision of holy ones they will assure you, and change your fear into
joy'' (43; Josh 5:13). Even the power to cast out demons is no ground of
boasting, ``for the working of signs is not ours but the Saviour's
work'' (38). Athanasius relates at the end of the \emph{Life} that Antony,
caught up in spirit, saw his soul led through the air by his conductors
while ``bitter and terrible beings'' sought to bar the way, and learned
``against what mighty opponents our wrestling is'' (65).

\sectionguard
\subsection{Cyril of Jerusalem}\label{sec:fb-cyril}

Cyril's \emph{Catecheses} are lectures to those preparing for baptism and
to the newly baptized. The angels appear in them as the catechumen meets
them in the creed and in the rites: as the witnesses of his baptism, as
creatures of the Son and servants of the Spirit, as the hosts of the
judgment, and as the choir whose hymn he joins at the altar. Because the
lectures teach the faith to beginners, their angelology is the Church's
common teaching in its plainest form.

\paragraph{The devil, once an archangel.} Cyril opens the matter of sin
with its first author. The devil ``sinned, not as having received
necessarily from nature the propensity to sin \dots\ but having been
created good, he has of his own free will become a devil, and received
that name from his action. For being an Archangel he was afterwards
called a devil from his slandering: from being a good servant of God he
has become rightly named Satan; for `Satan' is interpreted the
adversary'' (\emph{Cat.}\ II.4). The proof is Ezekiel's lament over the
king of Tyre, which Cyril reads of the devil: ``Thou wast perfect in thy
ways from the day of thy creation, until iniquity was found in thee''
(Ezek 28:15). ``Very rightly hath he said, were found in thee; for they
were not brought in from without, but thou didst thyself beget the
evil'' (II.4). The fall is final. The devil disguises himself as an
angel of light, ``not that he may reascend to where he was, for having
made his heart hard as an anvil, he has henceforth a will that cannot
repent'' (IV.1; 2 Cor 11:14). God rules over him and bears with him, not
from weakness but for two purposes, ``that he might disgrace himself the
more in his defeat, and that mankind might be crowned with victory'';
men are ``greatly honoured as having conquered him who was once an
Archangel'' (VIII.4).

\paragraph{Creatures of the Son.} Against those who ascribe the world to
angels, Cyril sets the creed's ``by whom all things were made'': ``For
whether visible or invisible, whether thrones or dominions, or anything
that is named, all things were made by Christ''; ``Silenced be they who
say that the world is the workmanship of Angels, who wish to steal away
the dignity of the Only-begotten'' (XI.21; Col 1:16). ``Christ made all
things, whether thou speak of Angels, or Archangels, of Dominions, or
Thrones'' (XI.23). The orders are real and ranked, and their very names
set a limit to curiosity. Cyril imagines the catechumen climbing through
the heavens to ask the angels how the Son was begotten; the angels of
the first heaven answer, ``We have above us beings greater and higher;
ask them,'' and the Thrones, Dominions, Principalities, and Powers
decline the question too, ``for they know it not'' (XI.11). ``Tell me
first, O most daring man, wherein does Throne differ from Dominion, and
then scrutinise what pertains to Christ. Tell me what is a Principality,
and what a Power, and what a Virtue, and what an Angel: and then search
out their Creator'' (XI.12). ``Be not ashamed to confess thine
ignorance, since thou sharest ignorance with Angels'' (XI.13).

\paragraph{Vision according to capacity.} The angels of the little ones
behold the Father's face (Matt 18:10), but, Cyril explains, ``the Angels
see God not as He is, but as far as they themselves are capable \dots\
The Angels therefore behold as much as they can bear, and Archangels as
much as they are able; and Thrones and Dominions more than the former,
but yet less than His worthiness: for with the Son the Holy Ghost alone
can rightly behold Him'' (VI.6). The measure of each order's vision is
the measure of its rank; the Only-begotten reveals the Father ``to each
according to his own capacity \dots\ through the Holy Ghost.'' Aquinas
draws the line between vision and comprehension: the blessed angels see
God's essence, but no created intellect comprehends it (\ST{I q.~12 a.~1 ad~1};
q.~12 a.~7).

\paragraph{Servants of the Spirit.} The name ``spirit'' is common to
many things: ``Thus an Angel is called spirit, our soul is called spirit,
and this wind which is blowing is called spirit \dots\ and a devil our
adversary is called spirit'' (XVI.13; Ps 103[104]:4). The Holy Spirit
is distinguished from them all as Lord. Cyril leads his hearers up
through the heavens in thought:

\begin{quote}
Ascend, I say, in imagination even unto the first heaven, and behold
there so many countless myriads of Angels. Mount up in thy thoughts, if
thou canst, yet higher; consider, I pray thee, the Archangels, consider
also the Spirits; consider the Virtues, consider the Principalities,
consider the Powers, consider the Thrones, consider the Dominions;---of
all these the Comforter is the Ruler from God, and the Teacher, and the
Sanctifier. Of Him Elias has need, and Elisseus, and Esaias, among men;
of Him Michael and Gabriel have need among Angels. Naught of things
created is equal in honour to Him. (XVI.23)
\end{quote}

\noindent The ascent is ordered: angels, archangels, spirits, virtues,
principalities, powers, thrones, dominions, with the Dominions at the
summit. The angels ``are sent forth to minister,'' but the Spirit
``searches even the deep things of God'' (XVI.23; Heb 1:14; 1 Cor 2:10).
The same lecture contrasts the coming of the Spirit with the assault of
the unclean spirit. The devil comes upon a soul ``like a wolf upon a
sheep, ravening for blood \dots\ His coming is most fierce; the sense of
it most oppressive; the mind becomes darkened'' (XVI.15). The Spirit's
coming is otherwise: ``First, His coming is gentle; the perception of
Him is fragrant; His burden most light; beams of light and knowledge
gleam forth before His coming. He comes with the bowels of a true
guardian'' (XVI.16). Cyril's contrast agrees with Antony's rule of
discernment.

\paragraph{The number of the angels and the judgment.} Cyril reckons the
multitude at the judgment: every nation now living, the barbarians, the
dead of a thousand years, all from Adam. ``Great indeed is the
multitude; but yet it is little, for the Angels are many more. They are
the ninety and nine sheep, but mankind is the single one'' (XV.24; Luke
15:4). He argues from the size of the heavens: ``The whole earth is but
as a point in the midst of the one heaven, and yet contains so great a
multitude; what a multitude must the heaven which encircles it contain?''
Daniel's ``thousand thousands'' are named ``not that the multitude is
only so great, but because the Prophet could not express more than
these'' (XV.24; Dan 7:10). Aquinas reasons in the same way from the
greatness of the heavens: as incorruptible bodies exceed the corruptible
in magnitude, so ``the immaterial substances as it were incomparably exceed
material substances as to multitude'' (\ST{I q.~50 a.~3}). At the
judgment the angel hosts encompass the
enemies of Christ so that they cannot flee, and gather the elect from the
four winds; ``though thou be in the field, the Angels shall take thee''
(XV.22--23). The virgin and the chaste will ``shine as an Angel''
(XV.23).

\paragraph{The angels in the rites.} The candidates are told that each
``is about to be presented to God before tens of thousands of the
Angelic Hosts'' (III.3), and that at their washing ``Angels shall dance
around you, and say, Who is this that cometh up in white array, leaning
upon her beloved?'' (III.16; Cant 8:5). After baptism Cyril explains the
anaphora. Following the thanksgiving, the priest and people make
mention of the whole creation,

\begin{quote}
of Angels, Archangels, Virtues, Dominions, Principalities, Powers,
Thrones; of the Cherubim with many faces: in effect repeating that call
of David's Magnify the Lord with me. We make mention also of the
Seraphim, whom Esaias in the Holy Spirit saw standing around the throne
of God, and with two of their wings veiling their face, and with twain
their feet, while with twain they did fly, crying Holy, Holy, Holy, is
the Lord of Sabaoth. For the reason of our reciting this confession of
God, delivered down to us from the Seraphim, is this, that so we may be
partakers with the hosts of the world above in their Hymn of praise.
(\emph{Mystagogical Catechesis} V.6; Ps 33[34]:4; Isa 6:2--3)
\end{quote}

\noindent Cyril names nine orders, from the Angels to the Seraphim, and
gives the reason for the Sanctus: the Church receives the hymn from the
Seraphim in order to share it. Aquinas explains the same place in the
Roman Mass: after the Preface ``the people devoutly praise Christ's
Godhead, saying with the angels: `Holy, Holy, Holy'\,'' (\ST{III q.~83
a.~4}); and the Catechism teaches that ``in her liturgy, the Church joins
with the angels to adore the thrice-holy God'' (CCC~335).

\sectionguard
\subsection{Basil of Caesarea}\label{sec:fb-basil}

Basil's teaching on the angels is spread through his defence of the
Holy Spirit, his homilies on the six days, his books against Eunomius,
and his homilies on the Psalms. Its centre is a doctrine of sanctification:
the angels are holy, not by nature, but by the Spirit.

\subsubsection{Holy by the Spirit}

In \emph{De Spiritu Sancto} Basil argues that the Spirit is inseparable
from the Father and the Son in every work, and he takes his first proof
from the first creation:

\begin{quote}
Moreover, from the things created at the beginning may be learnt the
fellowship of the Spirit with the Father and the Son. The pure,
intelligent, and supermundane powers are and are styled holy, because
they have their holiness of the grace given by the Holy Spirit.
(\emph{De Spiritu Sancto} 16.38)
\end{quote}

\noindent Moses tells only of the visible creation, but the believer may
reason from the seen to the unseen and ``glorify the Maker by whom all
things were made, visible and invisible, principalities and powers,
authorities, thrones, and dominions, and all other reasonable natures
whom we cannot name'' (16.38; Col 1:16). In that creation the three
Persons work as one: ``the ministering spirits subsist by the will of
the Father, are brought into being by the operation of the Son, and
perfected by the presence of the Spirit. Moreover, the perfection of
angels is sanctification and continuance in it'' (16.38). The Father,
the Word, and the Spirit are ``the Lord who gives the order, the Word who
creates, and the Spirit who confirms,'' and confirmation is ``the
perfecting according to holiness \dots\ firmness, unchangeableness, and
fixity in good'' (16.38; Ps 32[33]:6).

The holiness of the angels is therefore received, and received in
measure:

\begin{quote}
The powers of the heavens are not holy by nature; were it so there would
in this respect be no difference between them and the Holy Spirit. It is
in proportion to their relative excellence that they have their meed of
holiness from the Spirit. The branding-iron is conceived of together with
the fire; and yet the material and the fire are distinct. Thus too in
the case of the heavenly powers \dots\ their sanctification, being
external to their substance, superinduces their perfection through the
communion of the Spirit. They keep their rank by their abiding in the
good and true, and while they retain their freedom of will, never fall
away from their patient attendance on Him who is truly good. (16.38)
\end{quote}

\noindent The fall of some establishes the freedom of all: the wicked
spirits' ``fall establishes our statement of the freedom of the will of
the invisible powers; being, as they are, in a condition of equipoise
between virtue and vice, and on this account needing the succour of the
Spirit'' (16.38). Take away the Spirit, and ``the hosts of the angels are
disbanded, the dominions of archangels are destroyed, all is thrown into
confusion, and their life loses law, order, and distinctness.'' Gabriel
foretells by the Spirit's foreknowledge; the thrones and dominions live
their blessed life because they behold the Father's face, ``but to behold
it is impossible without the Spirit!'' Basil gives the image of a house
at night: withdraw the light, and gold is trodden on as though it were
iron. The heavenly order would lose its law like an army without its
general, ``or a chorus its harmony without the guidance of the
Coryph\ae us.'' ``How could the Seraphim cry `Holy, Holy, Holy,' were they
not taught by the Spirit how often true religion requires them to lift
their voice in this ascription of glory?'' The angels are ``perfect from
the moment of the creation,'' and in their creation ``there is \dots\ the
presence of the Holy Spirit, who confers on them the grace that flows from
Him for the completion and perfection of their essence'' (16.38).

Three determinations of the Summa answer to this chapter. Aquinas holds
it ``more probable, and more in keeping with the sayings of holy men,''
that the angels were created in sanctifying grace (\ST{I q.~62 a.~3});
that grace and glory were given them ``according to the degree of their
natural gifts'' (\ST{I q.~62 a.~6}), as Basil's ``in proportion to their
relative excellence''; and that the blessed angels cannot sin (\ST{I
q.~62 a.~8}), which is Basil's ``fixity in good.''

Basil states the angels' nature with a reserve that Aquinas does not
share: ``their substance is, peradventure, an aerial spirit, or an
immaterial fire \dots\ wherefore they exist in space and become visible,
and appear in their proper bodily form to them that are worthy'' (16.38;
Ps 103[104]:4). Elsewhere he sets the angels in place against the
Spirit's omnipresence: ``The angel who stood by Cornelius was not at one
and the same moment with Philip''; the Spirit alone fills the world
(23.54). Aquinas determines that the angels are altogether incorporeal
and assume bodies only for our sake (\ST{I q.~50 a.~1}; q.~51 aa.~1--2),
and that an angel is in a place by applying his power there, and in one
place only at a time (\ST{I q.~52 aa.~1--2}).

\paragraph{Witnesses, tutors, and guardians.} Why does Paul charge
Timothy ``before God and Christ Jesus and the elect angels'' (1 Tim
5:21)? Not because the angels are glorified with the Father and the
Son. ``The Spirit is called on as Lord of life, and the angels as allies
of their fellow-slaves and faithful witnesses of the truth.'' The saints
deliver God's commands before witnesses, and the angels will be present
when the Lord comes to judge (\emph{De Spiritu Sancto} 13.29; Luke 12:8).
``And so the apostle, knowing the angels to be set over men as tutors
and guardians, calls them to witness'' (13.30).

\subsubsection{The intelligible creation}

In the first homily on the six days Basil speaks of a creation before the
visible world:

\begin{quote}
It appears, indeed, that even before this world an order of things
existed of which our mind can form an idea, but of which we can say
nothing, because it is too lofty a subject for men who are but beginners
and are still babes in knowledge. \dots\ The Creator and Demiurge of the
universe perfected His works in it, spiritual light for the happiness of
all who love the Lord, intellectual and invisible natures, all the
orderly arrangement of pure intelligences who are beyond the reach of our
mind and of whom we cannot even discover the names. They fill the essence
of this invisible world, as Paul teaches us. ``For by him were all things
created that are in heaven, and that are in earth, visible and invisible
whether they be thrones or dominions or principalities or powers'' or
virtues or hosts of angels or the dignities of archangels.
(\emph{Hexaemeron} I.5)
\end{quote}

\noindent To this invisible world God added the visible one, ``both a
school and training place where the souls of men should be taught and a
home for beings destined to be born and to die'' (I.5). The intelligible
creation was not in darkness: ``The orders of angels, the heavenly hosts,
all intellectual natures named or unnamed, all the ministering spirits,
did not live in darkness, but enjoyed a condition fitted for them in
light and spiritual joy'' (II.5). When Basil later speaks of the vine as
an image of the soul, he names the angels among its supports: God has
surrounded human souls ``with the authority of His precepts and a guard of
angels. `The angel of the Lord encampeth round about them that fear
him'\,'' (V.6; Ps 33[34]:8). Aquinas records that the Greek Fathers held
the angels created before the visible world, and judges the contrary more
probable without deeming the Greek opinion erroneous (\ST{I q.~61 a.~3}
and ad~1; see Section~\ref{sec:q61}).

\subsubsection{The angel beside each believer}

The third book against Eunomius answers the claim that the Spirit, being
third in order, must be of another nature. Basil's reply turns on the
angels, who differ in rank and are one in nature:

\begin{quote}
All the angels, as they share one name, so also have altogether the same
nature with one another; yet some of them preside over nations, and
others accompany each one of the faithful. And by as much as a whole
nation is more to be esteemed than one man, by so much must the dignity
of the angel who rules a nation be greater than that of the angel
entrusted with the protection of each single person. That an angel is
with each of the faithful, as a kind of tutor and shepherd directing his
life, no one will gainsay who remembers the words of the Lord: ``See that
you despise not one of these little ones: for their angels always see
the face of my Father.''\footnote{Rendered for this edition from the
Greek of Migne, PG~29, 656B--657A; so also the sentence of III.2 below.} (\emph{Adversus Eunomium} III.1)
\end{quote}

\noindent The angel at each believer's side is a \gk{paidag\=ogos}, the
tutor who leads a child, and a \gk{nomeus}, a shepherd. The Catechism
quotes the sentence in its teaching that human life ``from infancy to
death'' is surrounded by the angels' watchful care: ``Beside each believer
stands an angel as protector and shepherd leading him to life'' (CCC~336).
Basil gathers his proofs. The Psalmist says, ``The angel of the Lord shall
encamp round about them that fear him'' (Ps 33[34]:8), and Jacob, in the
Septuagint that Basil quotes, names ``the angel who delivered me from my
youth'' (Gen 48:16). That angels
preside over whole nations ``Moses teaches us through the song,'' where
the Most High divided the nations and ``set the bounds of the nations
according to the number of his angels'' (Deut 32:8, which Basil reads in
the Septuagint; the Vulgate has ``the children of Israel''). Daniel heard
the angel say that ``the prince of the kingdom of the Persians'' withstood
him and that ``Michael, one of the chief princes, came to help'' him, and
that ``the prince of the Greeks'' was coming (Dan 10:13, 20). A
\gk{archistrat\=egos} of the Lord's host appeared to Joshua at the
Jordan; the Lord spoke of ``more than twelve legions of angels''; and the
chief captain of angels ranked in legions ``is plainly a ruler'' (Josh
5:14; Matt 26:53). The conclusion concerns the Spirit, but it states
Basil's angelology exactly: ``as among the angels one is a ruler and
another a subject, yet all are angels by nature, the difference is in
dignity and the communion is in nature'' (III.2).

Aquinas gives the same reason for the rank of the Principalities, who
preside ``over the government of peoples and kingdoms,'' because ``the
good of a nation is more divine than the good of one man,'' and cites the
prince of the Persians (\ST{I q.~108 a.~6}; Dan 10:13). The individual
guardian he places in the lowest order, ``whose place it is \dots\ to
announce the `lesser things'\,'' (\ST{I q.~113 a.~3}).

\subsubsection{The homilies on the Psalms}

Basil's homily on Psalm 33 gives the guardian angel a pastoral warmth
that the treatise against Eunomius does not need. On the verse ``The
angel of the Lord shall encamp round about them that fear him: and shall
deliver them,'' he says:

\begin{quote}
An angel attends upon every one who has believed in the Lord, unless we
ourselves drive him away by evil works. For as smoke puts bees to flight
and a foul smell drives out doves, so the lamentable and foul-smelling
sin drives away the angel who guards our life. If you have in your soul
works worthy of an angel's guarding, and a mind rich in the
contemplations of truth dwells in you through the wealth of the precious
works of virtue, God of necessity sets watchers and guardians beside you
and walls you round with a guard of angels.\footnote{Rendered for this
edition from the Greek of PG~29 (\emph{Homilia in Psalmum} 33.5, on v.~8).}
(\emph{Hom.\ in Ps.}\ 33.5)
\end{quote}

\noindent The psalm says the angel will ``encamp,'' and Basil marvels
that one angel is likened to a whole army and a crowded camp: ``because
of the greatness of him who guards you the Lord bestows on you a camp,
and because of the angel's strength he walls you round, as it were, on
every side with the safety that comes from him.'' As the walls of a city
keep out the enemy on every side, ``so too the angel both fortifies you
in front and guards your rear, and leaves neither side unguarded''
(33.5). The same homily teaches the fear of God by the judgment. Before
the dread tribunal ``there stand beside him who has done many evil things
in his life certain fearful and gloomy angels, looking fire, breathing
fire, because of the bitterness of their will, their faces like night
for their gloom and hatred of men'' (33.8; Ps 33[34]:12). And on the
verse that sets the eyes of the Lord on the just and his face against
evildoers, Basil recalls Matt 18:10: man cannot see God's face and live,
``but the angels of the little ones in the Church always behold the face
of our Father who is in heaven.'' Our flesh is a veil; ``the angels,
having no such covering like our flesh, are hindered by nothing from
gazing without ceasing on the face of the glory of God,'' and we too
shall see face to face when we become sons of the resurrection
(\emph{Hom.\ in Ps.}\ 33, on vv.~16--17).

In the homily on Psalm 48 Basil reads ``the just shall have dominion
over them in the morning'' (Ps 48[49]:15) of the angels: God ``handed
them over to the upright, that is, to his holy angels, to shepherd them.''
``For to each of the faithful an angel is joined, worthy to behold the
Father who is in heaven'' (\emph{Hom.\ in Ps.}\ 48.9). Basil's teaching
that sin drives the angel off is answered in the Summa by the rule that
the guardian ``never forsakes a man entirely, but sometimes he leaves him
in some particular,'' by not preventing trouble or even sin, ``according
to the ordering of Divine judgments'' (\ST{I q.~113 a.~6}).

\sectionguard
\subsection{Gregory of Nazianzus}\label{sec:fb-nazianzen}

Damascene, as Aquinas quotes him, calls Gregory ``the Theologian,'' and
Aquinas, through Jerome, cites his authority as one against which no one
has raised objection (\ST{I q.~61 a.~3}, obj.~1 and corpus). His teaching on the angels is
woven into his orations: the second theological oration, the festal
sermons on the Nativity, the Lights, and Pentecost, and his farewell to
Constantinople. He speaks of the angels as lights, the first overflow of
the Light that is God, and as the choir that circles the First Cause.

\paragraph{The natures and the names.} In the second theological oration
Gregory leads his hearers from the visible world to the invisible one.
Taking the Tabernacle as a figure of the whole creation, he asks whether
we shall ``pass the first veil, and stepping beyond the realm of sense,
\dots\ look into the Holy Place, the Intellectual and Celestial
creation.'' Even this we cannot see incorporeally, ``since it is
called---or is---Fire and Spirit.'' He reads the psalm verse of the angels'
nature:

\begin{quote}
The Angel then is called spirit and fire; Spirit, as being a creature of
the intellectual sphere; Fire, as being of a purifying nature; for I know
that the same names belong to the First Nature. But, relatively to us at
least, we must reckon the Angelic Nature incorporeal, or at any rate as
nearly so as possible. (\emph{Oratio} 28.31; Ps 103[104]:4)
\end{quote}

\noindent Then, as far as the mind can go, he names what we know:

\begin{quote}
\dots\ we know there are Angels and Archangels, Thrones, Dominions,
Princedoms, Powers, Splendours, Ascents, Intelligent Powers or
Intelligencies, pure natures and unalloyed, immovable to evil, or
scarcely movable; ever circling in chorus round the First Cause (or how
should we sing their praises?) illuminated thence with the purest
Illumination, or one in one degree and one in another, proportionally to
their nature and rank \dots\ so conformed to beauty and moulded that they
become secondary Lights, and can enlighten others by the overflowings
and largesses of the First Light. (28.31)
\end{quote}

\noindent The list joins six names that Scripture gives to three names drawn
from the angels' activity: Splendours,
Ascents, Intelligences. Gregory does not rank them. He continues with
their offices: they are ``Ministrants of God's Will, strong with both
inborn and imparted strength, traversing all space, readily present to
all at any place through their zeal for ministry and the agility of their
nature \dots\ different individuals of them embracing different parts of
the world, or appointed over different districts of the Universe, as He
knoweth who ordered and distributed it all'' (28.31). They are ``Hymners
of the Majesty of the Godhead, eternally contemplating the Eternal
Glory, not that God may thereby gain an increase of glory \dots\ but that
there may never be a cessation of blessings to these first natures after
God.'' His purpose is humility: ``even the secondary natures surpass the
power of our intellect; much more then the First'' (28.31).

\paragraph{The first creation and the fall of Lucifer.} In the oration
on the Nativity Gregory sets the angels first in the order of creation.
God's goodness could not rest in contemplating itself; ``Good must be
poured out and go forth beyond Itself to multiply the objects of Its
beneficence.''

\begin{quote}
He first conceived the Heavenly and Angelic Powers. And this conception
was a work fulfilled by His Word, and perfected by His Spirit. And so the
secondary Splendours came into being, as the Ministers of the Primary
Splendour; whether we are to conceive of them as intelligent Spirits, or
as Fire of an immaterial and incorruptible kind, or as some other nature
approaching this as near as may be. I should like to say that they were
incapable of movement in the direction of evil, and susceptible only of
the movement of good, as being about God, and illumined with the first
rays from God \dots\ but I am obliged to stop short of saying that, and to
conceive and speak of them only as difficult to move because of him, who
for his splendour was called Lucifer, but became and is called Darkness
through his pride; and the apostate hosts who are subject to him,
creators of evil by their revolt against good and our inciters.
(\emph{Or.}\ 38.9)
\end{quote}

\noindent The fall of Lucifer through pride is Gregory's reason for
calling the angelic nature difficult to move to evil rather than immovable.
``Then when His first creation was in good order, He conceives a second
world, material and visible'' (38.10). The intelligible natures are
``akin to Deity''; the sensible are ``utterly alien to It'' (38.10). Man
is made last, from both worlds: ``a new Angel, a mingled worshipper,
fully initiated into the visible creation, but only partially into the
intellectual'' (38.11). Aquinas reports this oration through Damascene,
``He first of all devised the angelic and heavenly powers, and the
devising was the making thereof,'' and on its account refuses to call
the Greek opinion erroneous (\ST{I q.~61 a.~3} obj.~1 and corpus). On the
first sin Aquinas agrees with Gregory that it was pride (\ST{I q.~63
a.~2}).

\paragraph{Lights from the Light.} The oration on the Holy Lights places
the angel in a descending order of lights. ``God is Light,'' the highest
and unapproachable.

\begin{quote}
A second Light is the Angel, a kind of outflow or communication of that
first Light, drawing its illumination from its inclination and obedience
thereto; and I know not whether its illumination is distributed according
to the order of its state, or whether its order is due to the respective
measures of its illumination. A third Light is man. (\emph{Or.}\ 40.5)
\end{quote}

\noindent Gregory leaves open whether rank follows illumination or
illumination rank. Aquinas answers that grace and glory were given the
angels according to the degree of their natural gifts (\ST{I q.~62
a.~6}), so that the orders are ``adequately distinguished by the gifts of
grace, but dispositively by natural gifts'' (\ST{I q.~108 a.~4}).
Sinlessness belongs to God; ``I venture to say also that it belongs to
the Angelic nature too; or at least, I would affirm that nature to be very
nearly sinless, because of its nearness to God'' (40.7). In the second
Easter oration the angels are ``the first luminous nature, next to the
first nature of all, because springing directly from it,'' and ``scarcely
Angels themselves could offer gifts worthy of its rank, those first and
intellectual and pure beings, who are also eye-witnesses of the Glory
That is on high'' (45.2).

\paragraph{The Spirit in the angels.} With Basil, Gregory ascribes the
angels' firmness to the Holy Spirit. At Pentecost he preaches: ``He
wrought first in the heavenly and angelic powers, and such as are first
after God and around God. For from no other source flows their perfection
and their brightness, and the difficulty or impossibility of moving them
to sin, but from the Holy Ghost'' (\emph{Or.}\ 41.11). The Spirit is
``penetrating all spirits that are intelligent, pure, most subtle (the
Angel Hosts I think)'' (31.29). Yet the angels are creatures and share in
some measure the mutability of creatures: ``perhaps the same may be said
of the Angels and the whole of that superior nature which is second to
the Trinity alone; although they are simple in some measure and more fixed
in good, owing to their nearness to the highest Good'' (31.15).

\paragraph{The angels of the churches and the feast.} Gregory holds that
each church has its angel. In his last farewell to Constantinople, spoken
before the bishops of the Council of 381, he hears the Lord call to
``the presiding Angels, for I believe, as John teaches me in his
Revelation, that each Church has its guardian,'' bidding them prepare the
way of his people (\emph{Or.}\ 42.9; Apoc 1:20; 2:1; Isa 62:10). Taking leave of
the city, he closes the same oration: ``Last of all, and most of all, I will
cry,---farewell ye Angels, guardians of this church, and of my presence and
pilgrimage, since our affairs are in the hands of God'' (42.27). The
priest ``is to take his stand with Angels, and give glory with
Archangels, and cause the sacrifice to ascend to the altar on high''
(2.73). At the Nativity the Church keeps feast with the angels: ``With
Shepherds glorify Him; with Angels join in chorus; with Archangels sing
hymns. Let this Festival be common to the powers in heaven and to the
powers upon earth. For I am persuaded that the Heavenly Hosts join in our
exultation and keep high Festival with us to-day \dots\ because they love
men, and they love God'' (38.17).

\sectionguard
\subsection{Gregory of Nyssa}\label{sec:fb-nyssa}

Basil's younger brother gives the angels a place in his account of the
whole creation, of man's fall, and of Christ's victory. Three teachings
stand out in him: that the devil fell through envy of man; that the
multiplication of men would have followed the angels' manner had man not
sinned; and that an angel and a demon stand at each man's side.

\paragraph{The fall of the power set over the earth.} In the
\emph{Catechetical Oration} Gregory explains the origin of evil by ``an
argument such as the following we have received by tradition from the
Fathers.'' God joined the intellectual and the sensible worlds, so that
nothing in creation should be without its share in the higher, and man is
the meeting of both. The angelic powers came first, each with its
charge:

\begin{quote}
Since, then, the intellectual nature had a previous existence, and to
each of the angelic powers a certain operation was assigned, for the
organization of the whole, by the authority that presides over all
things, there was a certain power ordained to hold together and sway the
earthly region, constituted for this purpose by the power that
administers the Universe. Upon that there was fashioned that thing
moulded of earth, an ``image'' copied from the superior Power. Now this
living being was man. \dots\ He to whom the administration of the earth
has been consigned takes it ill and thinks it not to be borne, if, of
that nature which has been subjected to him, any being shall be exhibited
bearing likeness to his transcendent dignity. (\emph{Oratio catechetica}
6)
\end{quote}

\noindent Evil has no being of its own; it is ``the deprivation of
goodness, just as a shadow which supervenes at the passage of the solar
ray.'' The uncreated alone is unchangeable; the created power ``could
choose by a spontaneous movement whatever he liked,'' and ``when he had
closed his eyes to the good \dots\ like one who in the sunshine lets his
eyelids down upon his eyes and sees only darkness,'' he became cognisant
of the contrary to goodness. ``Now this is Envy'' (6). Once turned, he
fell ``like a rock, torn asunder from a mountain ridge, which is driven
down headlong by its own weight,'' and made his intellect the instrument
of vice, deceiving man ``to become his own murderer with his own hands''
(6). Aquinas places pride first in the angel's sin and envy after it, the
devil grieving ``over man's good, and also over the Divine excellence''
(\ST{I q.~63 a.~2}). He reports the like teaching of Damascene, ``that
the highest of those who sinned was set over the terrestrial order''
(\ST{I q.~63 a.~7}; see Section~\ref{sec:q63}).

Christ's victory repays the deceiver in kind. The devil could not bear
the unveiled presence of God, so ``the Deity was hidden under the veil of
our nature, that so, as with ravenous fish, the hook of the Deity might
be gulped down along with the bait of flesh, and thus, life being
introduced into the house of death, and light shining in darkness, that
which is diametrically opposed to light and life might vanish'' (24). ``He
who first deceived man by the bait of sensual pleasure is himself deceived
by the presentment of the human form'' (26). Gregory speaks of that
healing as reaching even ``him, too, who had wrought our ruin'' (26); the
Church confesses the punishment of the demons to be everlasting, and the
synod of Constantinople in 543 condemned the teaching that it will end or
that the demons will be restored (DS~411; Lateran IV, DS~801).

\paragraph{How the angels increase.} Gregory asks why man was made male
and female, if the resurrection restores him to the life of the angels,
who do not marry, as the Lord answered the Sadducees (Luke 20:34--36). The
armies of the angels ``are in countless myriads'' without marriage;

\begin{quote}
so, in the same way, if there had not come upon us as the result of sin a
change for the worse, and removal from equality with the angels, neither
should we have needed marriage that we might multiply; but whatever the
mode of increase in the angelic nature is (unspeakable and inconceivable
by human conjectures, except that it assuredly exists), it would have
operated also in the case of men, who were ``made a little lower than the
angels,'' to increase mankind to the measure determined by its Maker.
(\emph{De hominis opificio} 17.2)
\end{quote}

\noindent God foresaw the fall and gave man a manner of increase suited
to fallen nature, ``in order that the multitude of human souls might not
be cut short by its fall from that mode by which the angels were
increased and multiplied'' (17.4). Gregory's question has its own answer:
the angels ``exist in countless myriads, being one essence, and at the
same time numerically many'' (17.3). Aquinas, for his part, holds each
angel a species of its own and their number exceedingly great, ``far
beyond all material multitude'' (\ST{I q.~50 aa.~3--4}; see
Section~\ref{sec:q50}).

\paragraph{The hundredth sheep.} Against Eunomius Gregory expounds Heb
1:6, ``when he bringeth in the first begotten into the world,'' of
Christ's return in glory, when ``the company of all the angels worships
Him Who comes in the name of the First-begotten, in their rejoicing over
the restoration of men.'' The angels rejoice over each sinner rescued, and
judge our perdition their own loss; ``when the sheep is brought safe to
the hundred above, (and we surely---humanity that is to say---are that
sheep which the Good Shepherd saved by becoming the first-begotten,) then
especially will they offer, in their intense thanksgiving on our behalf,
their worship to God'' (\emph{Contra Eunomium} IV.3, in the NPNF
division; Luke 15:4--7).

\paragraph{The angel and the demon at each man's side.} In the
\emph{Life of Moses} Gregory reads the meeting of Moses and Aaron as the
soul's meeting with its angel. He introduces the teaching as a tradition:

\begin{quote}
There is a teaching, which has its credibility from the tradition of the
fathers, that says that when our nature fell into sin, God did not
disregard our fall and leave it without his providence, but set beside
the life of each one, as an ally, an angel of those who have received the
bodiless nature; while on the other side the corrupter of our nature
contrived the like, by some wicked and evil-working demon who does harm
to the life of man. And man, standing between the two, each of his
companions having an aim opposed to the other's, himself makes the one
prevail over the other.\footnote{Rendered for this edition from the Greek
of Migne, PG~44, 337D--340A.} (\emph{De vita Moysis} II)
\end{quote}

\noindent The good angel sets before the reason the goods of virtue that
hope sees; the other, the material pleasures that enslave the senses. When
a man turns from the bait of evil and sets his soul like a mirror toward
the hope of good, ``then the alliance of the brother meets him and stands
by him. For the angel is in a certain way a brother to the rational and
intellectual part of man's soul'' (PG~44, 340). It is not in doubt, he
says, that the angel is akin to the soul in its intelligible and bodiless
being, was made before us, and stands by those who wrestle with their
adversaries. The Church teaches that an angel watches over each
human life (CCC~336), and Aquinas teaches the assault of the demons as
permitted by God and ordered by him to man's good (\ST{I q.~114 a.~1}).

\sectionguard
\subsection{Ephrem the Syrian}\label{sec:fb-ephrem}

Ephrem writes in Syriac and in verse, and his angelology is sung. His
usual name for the angels is ``the Watchers,'' the wakeful ones; he
delights to set them beside the ``vigil-keepers'' of the Church, who watch
through the night of the Nativity. Quotations are from the hymns in the
thirteenth volume of the second series of the \emph{Nicene and Post-Nicene
Fathers}.

\paragraph{The Watchers at the Nativity.} ``Joyous were to-day the
Watchers, that the Wakeful came to wake us! Who would pass this night in
slumber, in which all the world was watching?'' (\emph{Hymns on the
Nativity} I). The descent of the angels to the shepherds joins heaven to
the Church's vigil:

\begin{quote}
To-day the angels, and the archangels,---descended to sing---a new song on
earth.---In this mystery they descend, and rejoice with the
vigil-keepers. \dots\ For this is the night that joined, the Watchers on
high with the vigil-keepers.---The Watcher came to make watchers in the
midst of creation.---Lo! the vigil-keepers are made comrades with the
Watchers:---the singers of praise are made, companions of the Seraphs.
(\emph{Nat.}\ XIV.3--4)
\end{quote}

\noindent Christ is ``the Watcher'' who came to make men watchers.
Gabriel, ``chief of Angels, called Him `My Lord' \dots\ to teach that He
was his Lord, not his fellow.---Gabriel had with him, Michael as fellow:---the
Son is Lord of the servants'' (XIV.23). The thousands who stand before
the throne cannot search him out, ``for all of them stand, in silence to
serve'' (XIV.21; Dan 7:10). One angel announced the conception, but a
multitude the birth, ``for greater is the joy of the Birth than the
Conception'' (XV.36--37). ``It was not Seraphim He sent us, nor yet did
Cherubim come down among us;---there came not Watchers or Ministers, but
the Firstborn to Whom they minister'' (XVI.8).

\paragraph{The Watchers at the font.} In the hymns for the Epiphany the
angels rejoice over the newly baptized. The lowly man who has put on
Christ in the water ``is great in the sight of the Watchers'' (\emph{Hymns
for the Epiphany} IV.10). ``The Angels and the Watchers rejoice---over that
which is born of the Spirit and of water \dots\ The Seraphins who sing
`Holy' rejoice,---that they who are made holy have been increased'' (VI.7).
``For lo! the Angels rejoice---over one sinner if he repent:---how much more
do they now rejoice---that in all churches and congregations,---lo! Baptism
is bringing forth---the heavenly from the earthly!'' (VI.8; Luke 15:10).
Baptism itself becomes a guard like the host about Elisha: ``Elisha was
the equal of the Watchers \dots\ The camp of the Watchers was round about
him;---thus let Baptism be unto you,---a camp of guardians'' (VIII.19; 4
Kgs 6:17). Another hymn of the series closes with the whole Church,
above and below: ``From
every mouth with one consent,---of those beneath and those above,---Watchers,
Cherubin, and Seraphin,---the baptized, the sealed, and the hearers,---let
each of us cry aloud and say,---`Glory to the Lord of our feasts!'\,''
(VI.20).

\paragraph{The Watchers in Hades.} In the Nisibene hymns Death speaks
of the Resurrection. The angels who sat in the tomb, one at the head and
one at the feet, are to him an invasion of his realm: ``Death has seen
the Watchers in Hell; the immortal instead of the mortal; and he said
Confusion has entered our abode \dots\ That the dead have come forth out
of Hell, and the Watchers that die not have entered therein'' (\emph{Carmina
Nisibena} XXXVI.15; John 20:12).

\sectionguard
\subsection{John Chrysostom}\label{sec:fb-chrysostom}

Chrysostom's homilies on Matthew, John, Acts, and the Epistles pause
again and again where the text names the angels. His angelology is
pastoral: the angels serve our salvation,
they are fellow-servants and tutors, they stand about the altar, and
their presence is a reason for sobriety and for joy.

\subsubsection{Ministering spirits}

The third homily on Hebrews expounds Heb 1:6--14. The Father ``bringeth in
the First-Begotten'' and commands the angels to worship him, and the
Apostle sets the angels' making against the Son's throne: ``Behold, the
greatest difference! that they are created, but He uncreated''
(\emph{Hom.\ in Heb.}\ 3.1). On the last verse Chrysostom turns to the
hearers:

\begin{quote}
Ver. 14. ``Are they not all ministering spirits, sent forth to minister
for them who shall be heirs of salvation?'' What marvel (saith he) if
they minister to the Son, when they minister even to our salvation? See
how he lifts up their minds, and shows the great honor which God has for
us, since He has assigned to Angels who are above us this ministration on
our behalf. \dots\ this is the office of Angels, to minister to God for
our salvation. So that it is an angelical work, to do all for the
salvation of the brethren: or rather it is the work of Christ Himself,
for He indeed saves as Lord, but they as servants. And we, though
servants are yet Angels' fellow-servants. (3.4)
\end{quote}

\noindent ``And yet the space between angels and men is great; nevertheless he
brings them down near to us, all but saying, For us they labor, for our
sake they run to and fro: on us, as one might say, they wait. This is
their ministry, for our sake to be sent every way'' (3.4). Chrysostom
runs through the examples: the shepherds, Mary, Joseph, the tomb, the
Ascension, Peter freed from prison, Philip, Cornelius, Paul; ``Dost thou
see that they minister to us on God's behalf, and that they minister to us
in the greatest matters?'' (3.4). The Son too was sent, but not as
servant; indeed he was not sent in the way the angels are, ``for He did
not pass from place to place, but took on Him flesh: whereas these change
their places, and leaving those in which they were before, so come to
others in which they were not.'' The Apostle's word is also comfort:
``What fear ye? Angels are ministering to us'' (3.4).

The glory of the angels makes the glory of our nature in Christ the
greater. On ``for nowhere doth he take hold of the angels'' (Heb 2:16)
Chrysostom exclaims: ``For in very deed it is a great and a wonderful
thing, and full of amazement that our flesh should sit on high, and be
adored by Angels and Archangels, by the Cherubim and the Seraphim. For
myself having oftentimes thought upon this, I am amazed at it'' (5.1).
Yet the angels do not comprehend God. On ``No man hath seen God at any
time'' (John 1:18) he teaches that ``what God really is, not only have not
the prophets seen, but not even angels nor archangels. If you ask them,
you shall not hear them answering anything concerning His Essence, but
sending up, `Glory to God in the Highest' \dots\ If you desire to learn
something from Cherubim or Seraphim, you shall hear the mystic song of
His Holiness'' (\emph{Hom.\ in Jo.}\ 15.1; Luke 2:14; Isa 6:3). Aquinas
cites this homily as the first objection to the vision of God's essence,
and answers that Chrysostom speaks ``of the vision of comprehension''
(\ST{I q.~12 a.~1} obj.~1 and ad~1).

\subsubsection{Principalities and powers}

In the homilies on Colossians Chrysostom meets teachers who would bring
men to God through the angels. On Col 1:16 he notes that Paul names the
thrones and dominions because their creation by Christ was in question:
``what is granted, he lets alone, but what is doubted he asserts'';
``whether'' and ``or'' embrace all things, ``but the Spirit is not amongst
the `powers'\,'' (\emph{Hom.\ in Col.}\ 3). The peace made by the blood of
the Cross (Col 1:20) was a peace with heaven as well as among men:

\begin{quote}
Why then place ye confidence in Angels? saith he. For so far are they
from bringing you near, that they were ever your enemies, except God
Himself had reconciled you with them. Why then run ye to them? Wouldest
thou know the hatred which the Angels had against us, how great it was;
and how averse to us they always were? They were sent to take vengeance
in the cases of the Israelites, of David, of the Sodomites, of the Valley
of weeping. Not so however now, but, on the contrary, they sang upon the
earth with exceeding joy. (\emph{Hom.\ in Col.}\ 3)
\end{quote}

\noindent ``He brought these first down hither, and then he took up man
to them; earth became heaven, because that heaven was about to receive
the things of earth'' (3). Human nature, once below the stones in its
idolatry, is now set ``higher in dominion than the Angels, and the
Archangels, and the thrones, and the dominions'' (5; Col 1:27). On the
``religion of angels'' (Col 2:18): ``There are some who maintain that we
must be brought near by Angels, not by Christ, that were too great a
thing for us''; such a man ``hath not seen Angels, and yet is affected as
though he had,'' and he has let go the Head: ``Why, letting go the Head,
dost thou cling to the members?'' (7).

The principalities of Col 2:15 are the fallen ones: ``He means the
diabolical powers.'' The devil, expecting to take Christ, ``lost even those
he had; and when That Body was nailed to the Cross, the dead arose. There
death received his wound, having met his death-stroke from a dead body''
(6). In Ephesians the powers against whom we wrestle are named as the
holy powers are named: ```Principalities, and powers,' he speaks of; just
as in the heavenly places there are `thrones and dominions, principalities
and powers'\,'' (\emph{Hom.\ in Eph.}\ 22; Eph 6:12; Col 1:16). The warning
follows: ``if `the principalities' are incorporeal, if they are `rulers of
the world,' if they are `the spiritual hosts of wickedness,' how, tell me,
canst thou live in self-indulgence?'' (22).

The mystery of the Church was hidden even from the angels. On Eph 3:10
Chrysostom has Paul answer his own question: ``art thou enlightening Angels
and Archangels and Principalities and Powers? I am, saith he \dots\ But
whence hath this been made manifest to the Angels? By the Church.'' ``When
we were taught it, then were they also by us'' (\emph{Hom.\ in Eph.}\ 7).
The angels knew that ``the Lord's portion was His people'' and that the
Gentiles would be called, ``but that it was to the same privileges as
Israel, yea, even to sit upon the throne of God, this, who would ever have
expected?'' (7; Deut 32:9). Aquinas teaches that the angels knew the
mystery of the Incarnation in general from the beginning and learned its
special conditions afterward (\ST{I q.~57 a.~5} and ad~1).

\subsubsection{The angel of each believer}

Three homilies give Chrysostom's teaching on the guardian angel. On Matt
18:10 he draws the conclusion broadly: ``Hence it is evident, that the
saints have angels, or even all men,'' and he cites 1 Cor 11:10 and the
Septuagint of Deut 32:8, ``according to the number of the angels of God.''
Of the text itself he says: ``here He is discoursing not of angels only,
but rather of angels that are greater than others. But when He saith,
`The face of my Father,' He means nothing else than their fuller
confidence, and their great honor'' (\emph{Hom.\ in Matt.}\ 59.4). Aquinas
cites the reading as an objection to his placing the guardians in the
lowest order, and answers that Chrysostom ``can be taken to mean the
highest in the lowest order of angels,'' adding that ``it is, however,
probable that the greater angels are deputed to keep those chosen by God
for the higher degree of glory'' (\ST{I q.~113 a.~3} obj.~1 and ad~1). On
Acts 12:15, where the disciples at Mary's door say of Peter, ``It is his
angel,'' Chrysostom states flatly: ``This is a truth, that each man has an
Angel'' (\emph{Hom.\ in Act.}\ 26). In the homily on Col 1:20 he ties the
guardian to the new order of the Church:

\begin{quote}
At first the Angels were according to the number of the nations; but now,
not to the number of the nations, but that of the believers. Whence is
this evident? Hear Christ saying, ``See that ye despise not one of these
little ones, for their Angels do always behold the face of My Father which
is in heaven.'' For each believer hath an Angel; since even from the
beginning, every one of those that were approved had his Angel, as Jacob
says, ``The Angel that feedeth me, and delivereth me from my youth.'' If
then we have Angels, let us be sober, as though we were in the presence of
tutors; for there is a demon present also. Therefore we pray, asking for
the Angel of peace. (\emph{Hom.\ in Col.}\ 3; Gen 48:15--16)
\end{quote}

\noindent Chrysostom speaks now of ``each man,'' now of ``each believer,''
as Basil does of ``each of the faithful.'' Aquinas records two opinions,
that the guardian is given at baptism and that he is given at birth, and
follows Jerome in the second: the guardian is given with the nature, not
with the grace of baptism (\ST{I q.~113 a.~5}; see
Section~\ref{sec:q113}). The same presence grounds a rule of modesty.
Paul bids the woman be veiled ``because of the angels'' (1 Cor 11:10), and
Chrysostom paraphrases: ```For although thou despise thine husband,' saith
he, `yet reverence the angels'\,'' (\emph{Hom.\ in 1 Cor.}\ 26).

\subsubsection{The angels about the altar}

Chrysostom returns again and again to the angels' presence at the
Eucharist. In the treatise \emph{On the Priesthood} he
sets the priest's office among the heavenly orders:

\begin{quote}
For the priestly office is indeed discharged on earth, but it ranks
amongst heavenly ordinances; and very naturally so: for neither man, nor
angel, nor archangel, nor any other created power, but the Paraclete
Himself, instituted this vocation, and persuaded men while still abiding
in the flesh to represent the ministry of angels. Wherefore the
consecrated priest ought to be as pure as if he were standing in the
heavens themselves in the midst of those powers. (\emph{De sacerdotio}
III.4)
\end{quote}

\noindent Priests ``have received an authority which God has not given
to angels or archangels,'' for to no angel was it said, ``Whatsoever ye
shall bind on earth shall be bound in Heaven'' (III.5; Matt 18:18). When
the priest invokes the Holy Spirit and offers ``the most dread
sacrifice,'' the angels are there:

\begin{quote}
At such a time angels stand by the Priest; and the whole sanctuary, and
the space round about the altar, is filled with the powers of heaven, in
honor of Him who lieth thereon. For this, indeed, is capable of being
proved from the very rites which are being then celebrated. I myself,
moreover, have heard some one once relate, that a certain aged, venerable
man, accustomed to see revelations, used to tell him, that he being
thought worthy of a vision of this kind, at such a time, saw, on a
sudden, so far as was possible for him, a multitude of angels, clothed in
shining robes, and encircling the altar, and bending down, as one might
see soldiers in the presence of their King, and for my part I believe it.
(VI.4)
\end{quote}

\noindent Another told him that when those who have received the
mysteries with a pure conscience are about to die, ``angels keep guard
over them for the sake of what they have received, and bear them hence''
(VI.4). The same teaching sounds in the homilies. ``Look, I entreat: a
royal table is set before you, Angels minister at that table, the King
Himself is there, and dost thou stand gaping?'' (\emph{Hom.\ in Eph.}\ 3).
It is ``a sacrifice, at which the very Angels tremble'' (3). In the litany
for the catechumens the deacon bids them, ``Pray ye, Catechumens, for the
angel of peace''; Chrysostom explains, ``for there is an angel that
punisheth \dots\ Wherefore we bid them ask for the angel of peace, teaching
them to seek that which is the bond of all good things, peace'' (\emph{Hom.\
in 2 Cor.}\ 2; Ps 77[78]:49). Aquinas, expounding the Roman Canon, speaks
in the same way of ``the angel standing by at the Divine mysteries'' who
presents the prayers of priest and people to God (\ST{III q.~83 a.~4
ad~9}; Apoc 8:4).

The angels will also stand at the judgment, and fear with us. Chrysostom
bids the fallen Theodore see the heavenly assembly of ``angels,
archangels, thrones, dominions, principalities, powers'' (\emph{Ad
Theodorum lapsum} I.11), and at the judgment ``even the angels themselves
are holden by much fear, and not angels only but also archangels and
thrones, and dominions, and principalities and authorities,'' because
their fellow-servants are called to account (I.12).

\subsubsection{The demons do not govern the world}

Against those who excused their sins by saying that the world is given
over to demons, Chrysostom preached three homilies on the devil's power.
The demons' treatment of the Gerasene swine shows what their rule would
be: ``Thus do Demons govern'' (\emph{De diabolo tentatore} I.6). Had God
entrusted the world to them, ``they would have confused and disturbed
everything''; the order of the sun, stars, and seasons shows the
providence that governs it (I.6). The devil ``does not overcome by force,
nor by tyranny, nor through compulsion, nor through violence,'' as Job's
flocks and the swine show, for he could touch neither without permission
from above (II.1). Why then is he not removed? Because ``if the
antagonist were taken away he who overcomes is thereby injured'': the
diligent would ``neither exhibit their own power, nor win crowns'' (II.1).
Nor is the devil evil by nature. ``Let the Devil be allowed to be exceeding
wicked, not by nature, but by choice and conviction''; his very names,
slanderer and wicked one, name acts and not an essence, and his wickedness
``was not indeed with him at the beginning, but afterwards came upon him;
wherefore he is called apostate.'' He is called wicked above all others
because, though wronged in nothing, ``when he saw mankind had in honour,
he straightway envied him his good'' (II.2).
Aquinas determines the same of the assault of the demons: the
assault is the demons' malice, its ordering is God's, ``Who knows how to
make orderly use of evil by ordering it to good'' (\ST{I q.~114 a.~1}).

\sectionguard
\subsection{The Nicene Fathers and Their Reception}\label{sec:fb-nicene-reception}

The table gathers the teaching of the Greek and Syriac Fathers from
Athanasius to Chrysostom at the loci expounded above, with the place
where Aquinas or the Church receives it. A dash marks a teaching that
stands without a single receiving locus.

\begin{fourstable}{Witness}{Locus}{Teaching}{Received at}
Athanasius & \emph{C.\ Ar.}\ II.19--21, 26--29 & The orders are creatures; angels do not create; God needs no created mediator in order to create & \ST{I q.~45 a.~5}; Lateran IV (DS~800)\\
Athanasius & \emph{C.\ Ar.}\ I.55--62; II.23 & The Son is ``better'' than the angels by nature, not by degree; the angels worship him; worship belongs to God alone & CCC 331, 333\\
Athanasius & \emph{C.\ Ar.}\ III.12--14 & The angels minister at God's command; in the angel of the bush God spoke & \ST{I q.~51 a.~2 ad~1}\\
Athanasius & \emph{C.\ Ar.}\ III.10; \emph{De inc.}\ 25 & The good angels remain by willing what God wills; the fallen hold the lower air, cleared by the Cross & \ST{I q.~64 a.~4}\\
Athanasius & \emph{V.\ Ant.}\ 22, 28--30, 35, 41 & The demons were made good and fell; since Christ they are weak, act by permission, and flee the Cross and Christ's name & Lateran IV (DS~800); \ST{I q.~114 a.~1}\\
Athanasius & \emph{V.\ Ant.}\ 31--33 & The demons know the future only by conjecture, as physicians do & \ST{I q.~57 a.~3}\\
Athanasius & \emph{V.\ Ant.}\ 35--37, 43 & A holy visitation brings calm and joy; a demonic one, tumult and fear & ---\\
Cyril of Jerusalem & \emph{Cat.}\ II.4; IV.1; VIII.4 & The devil, once an archangel, fell by free will; his will cannot repent & \ST{I q.~63 aa.~4--5}; \ST{I q.~64 a.~2}\\
Cyril & \emph{Cat.}\ XI.21--23 & Christ made all the orders; the world is not the work of angels & CCC 331; \ST{I q.~45 a.~5}\\
Cyril & \emph{Cat.}\ VI.6; XI.11--13 & Each order sees God according to its capacity; none knows the Son's generation & \ST{I q.~12 aa.~1, 7}\\
Cyril & \emph{Cat.}\ XVI.23 & Angels, archangels, spirits, virtues, principalities, powers, thrones, dominions: the Spirit is their Ruler and Sanctifier & \ST{I q.~62 a.~3}\\
Cyril & \emph{Cat.}\ XV.24 & The angels are the ninety-nine sheep; their number exceeds all mankind & \ST{I q.~50 a.~3}\\
Cyril & \emph{Myst.}\ V.6 & Nine orders named in the anaphora; the Church sings the Seraphim's hymn to share it & \ST{III q.~83 a.~4}; CCC 335\\
Basil & \emph{De Sp.\ S.}\ 16.38 & The powers are holy by the Spirit, not by nature, in measure of their excellence; perfected at creation & \ST{I q.~62 aa.~3, 6, 8}\\
Basil & \emph{De Sp.\ S.}\ 16.38; 23.54 & The angels' substance is perhaps aerial spirit or immaterial fire; they are in one place & Answered: \ST{I q.~50 a.~1}; \ST{I q.~52 a.~2}\\
Basil & \emph{Hex.}\ I.5; II.5 & The intelligible creation and its orders existed before the visible world, in light & \ST{I q.~61 a.~3} and ad~1\\
Basil & \emph{Adv.\ Eun.}\ III.1 & An angel is beside each believer as tutor and shepherd; greater angels rule nations & CCC 336; \ST{I q.~108 a.~6}; \ST{I q.~113 a.~2}\\
Basil & \emph{Hom.\ in Ps.}\ 33.5; 48.9 & An angel attends every believer; sin drives him off; a guard of angels walls the just & \ST{I q.~113 aa.~4, 6}\\
Gregory Nazianzen & \emph{Or.}\ 28.31 & Angels, archangels, thrones, dominions, princedoms, powers, splendours, ascents, intelligences; ministers and hymners; set over regions & \ST{I q.~108 a.~5}; \ST{I q.~110 a.~1}\\
Gregory Nazianzen & \emph{Or.}\ 38.9--11 & God first conceived the angelic powers; Lucifer fell through pride; man is a ``new Angel'' & \ST{I q.~61 a.~3}; \ST{I q.~63 a.~2}\\
Gregory Nazianzen & \emph{Or.}\ 40.5, 7; 41.11 & The angel is a second light; its firmness in good is from the Spirit & \ST{I q.~62 aa.~6, 8}; \ST{I q.~106 a.~1}\\
Gregory Nazianzen & \emph{Or.}\ 42.9, 27 & Each church has its guardian angel & ---\\
Gregory of Nyssa & \emph{Or.\ cat.}\ 6 & The power set over the earth fell through envy of man & Pride first: \ST{I q.~63 a.~2}\\
Gregory of Nyssa & \emph{Or.\ cat.}\ 24, 26 & The deceiver is deceived by the hook of the Godhead & Restoration of the demons condemned: Constantinople 543 (DS~411)\\
Gregory of Nyssa & \emph{De hom.\ op.}\ 17 & The angels increase in their own manner, one essence and many in number & \ST{I q.~50 aa.~3--4}\\
Gregory of Nyssa & \emph{Vita Moysis} II & By tradition of the fathers, an angel and a demon attend each man & CCC 336; \ST{I q.~113 a.~2}; \ST{I q.~114 a.~1}\\
Ephrem & \emph{Nat.}\ XIV; XV--XVI & The Watchers sing at the Nativity; the Son is Lord of Gabriel and Michael & CCC 333\\
Ephrem & \emph{Epiph.}\ VI.7--8, 20; VIII.19 & The Watchers and Seraphim rejoice over the baptized; baptism is a camp of guardians & ---\\
Chrysostom & \emph{Hom.\ in Heb.}\ 3.4 & All angels are ministering spirits sent for our salvation; we are their fellow-servants & \ST{I q.~112 a.~1}; \ST{I q.~57 a.~5 ad~1}; CCC 331\\
Chrysostom & \emph{Hom.\ in Jo.}\ 15.1 & Not even the angels see what God is & \ST{I q.~12 a.~1} obj.~1 and ad~1\\
Chrysostom & \emph{Hom.\ in Col.}\ 3, 7; \emph{Hom.\ in Eph.}\ 7 & Angels were reconciled by the Cross; men are not brought to God by angels; the angels learn the mystery through the Church & \ST{I q.~57 a.~5}\\
Chrysostom & \emph{Hom.\ in Matt.}\ 59.4 & The angels of Matt 18:10 are greater angels; all men have angels & \ST{I q.~113 a.~3} obj.~1 and ad~1\\
Chrysostom & \emph{Hom.\ in Act.}\ 26; \emph{Hom.\ in Col.}\ 3 & Each man, each believer, has an angel; a demon too is present; pray for the angel of peace & \ST{I q.~113 aa.~2, 5}; CCC 336\\
Chrysostom & \emph{De sac.}\ III.4; VI.4; \emph{Hom.\ in Eph.}\ 3 & Angels fill the sanctuary at the sacrifice and minister at the table & \ST{III q.~83 a.~4 ad~9}; CCC 335\\
Chrysostom & \emph{De diab.\ tent.}\ I.6; II.1 & The demons do not govern the world; the devil overcomes by deceit, not force & \ST{I q.~114 a.~1}\\
\end{fourstable}
